\section{Efficient composition and certificate reuse}
\label{sec:efficiency}

The framework's efficiency claim is structural. We count evaluations
of a presented input, rational arithmetic operations, and distinct
certified observations. We do not equate primitive recursiveness with
practical polynomial time, nor claim a measured speedup for an
unimplemented proof compiler.

\subsection{Reverse sensitivity on a finite dependency graph}

Consider a finite directed acyclic graph, numbered so that every edge
\(u\to v\) has \(u<v\). Suppose a local estimate has been proved
for each node:
\[
 e_v\leq\eta_v+\sum_{u\to v}L_{vu}e_u,
 \qquad L_{vu}\in\mathbb Q_{\geq0}.
 \tag{L}\label{eq:localbudget}
\]
Here \(e_v\) is the error in the observation norm assigned to node
\(v\), and \(\eta_v\) is its local approximation error. A change of
norm is itself a node with a proved estimate. In particular, a
diagonal pullback cannot appear as an \(L^1\) edge with a fictitious
finite constant.

Let \(o\) be the output. Initialize \(w_o=1\) and all other weights
to zero. Visit nodes in reverse order and, for every edge \(u\to v\),
add \(w_vL_{vu}\) to \(w_u\).

\begin{theorem}[Shared error budget]\label{thm:dagbudget}
The reverse pass computes nonnegative rational weights satisfying
\[
 e_o\leq\sum_v w_v\eta_v.
 \tag{B}\label{eq:globalbudget}
\]
It uses one multiplication and one addition per edge. Every weight is
the sum of products of local sensitivities along all paths from that
node to the output; paths are not explicitly enumerated.
\end{theorem}
\begin{proof}
Start with \(e_o\). When the largest remaining error term \(w_ve_v\)
is processed, substitute \eqref{eq:localbudget} and distribute
\(w_v\). Its residual contribution is \(w_v\eta_v\), and each
predecessor contribution is \(w_vL_{vu}e_u\). No predecessor has yet
been processed, so contributions to an already shared node are added
before its one substitution. The invariant is an upper bound consisting
of processed residual terms plus unprocessed weighted error terms.
All coefficients are nonnegative, so substitution preserves the
inequality. At the end only residual terms remain. The path
description follows by the same recursion: prepend one edge to each
path already represented by \(w_v\).
\end{proof}

For a requested rational tolerance \(\varepsilon>0\), let
\(W=\sum_v w_v\) and choose every local error at most
\(\varepsilon/(1+W)\). Then the final error is strictly below
\(\varepsilon\). This uniform allocation is robust and needs no
optimization. More efficient allocations may use separate positive
budgets \(\beta_v\) with \(\sum_v\beta_v\leq\varepsilon\) and
require \(w_v\eta_v\leq\beta_v\). Zero weights need no demand.
All choices here use rational arithmetic and bounded dyadic searches.

\paragraph{Example with a shared input.}
If a node \(x\) is used in \(a=2x\) and \(b=3x\), followed by
\(o=a+b\), the reverse pass assigns weight \(5\) to the single
approximation of \(x\). It does not ask an expensive producer for two
independent values. With clipping, conditional expectation, or
integration, many corresponding constants are one. A finite moment
table similarly shares one step-function approximation across every
coefficient, as in \eqref{eq:moments}.

\subsection{The data that a certificate must retain}

A reusable finite certificate records the input identity, observation
norm, precision, common refinement, rational bounds, and actual trace
or fixed-program identifier used. Its proof obligations include
well-formed denominators and arrays, the local numerical inequality,
and any modulus instance that justified the demanded input precision.
A cache hit is valid only if it supplies at least that precision in the
same observation. Refinement certificates permit reuse across
different grid levels; they are not replaceable by matching filenames.

In a proof implementation, the fixed graph is compiled by substitution
and cut in the ordinary \(\Prf\,\mathrm{PRAAx}\) calculus. A
variable-sized graph requires a separately proved bounded checker or
replay invariant. The present reverse-pass theorem is not a universal
soundness theorem for arbitrary Lean proof objects, and it does not
reflect arbitrary checked formulas into \ipra.

\subsection{Concrete improvements over independent approximations}

\begin{itemize}
\item Dyadic refinement is delayed until the common level needed by
the current finite sum. Integral-preserving refinement identities are
proved once and reused across all moments, kernels, and formula nodes.
\item To obtain \(N\) moments of a step function with \(J\) cells,
update endpoint powers recursively. This needs \(O(JN)\) rational
operations after one input evaluation, rather than \(N\) calls to a
general integration procedure.
\item A density is first reconstructed from its energy certificates.
Every later integral uses its \(L^1\) name. Reconstructing a density
separately for each observable would repeat the hardest producer.
\item Stable topological output is requested with a scale interval.
Interval distance certificates are reused across all simplices.
The algorithm spends no work deciding a threshold equality whose
answer would not survive the input uncertainty.
\item A contraction gap gives an explicit finite iteration count.
No semidecision loop runs until an estimated residual appears small.
The residual is checked at a count already justified by the gap.
\end{itemize}

These improvements reduce duplicated demand and proof material. A
complete complexity theorem must additionally track rational bit sizes,
the representation of refinements, and the cost of the input producers.
All of those costs are bounded by fixed primitive-recursive programs
in the present setting. They may nevertheless grow rapidly; for
example a dyadic partition of dimension \(d\) and level \(m\) has
\(2^{dm}\) cells. No dimension-free computational claim is intended.
